\documentclass[11pt,a4paper]{article}
\usepackage[margin=21mm,headheight=15pt]{geometry}
\usepackage[T1]{fontenc}
\usepackage{lmodern,amsmath,amssymb,booktabs,tabularx,enumitem,xcolor,fancyhdr}
\usepackage[hidelinks]{hyperref}
\definecolor{accent}{RGB}{0,114,178}
\setlength{\parindent}{0pt}
\setlength{\parskip}{6pt}
\setlist{nosep,leftmargin=*}
\pagestyle{fancy}
\fancyhf{}
\fancyhead[L]{\small Econometrics 25117 | Instructor guide}
\fancyhead[R]{\small 2026--2027}
\fancyfoot[L]{\small Private teaching notes}
\fancyfoot[R]{\thepage}
\newcommand{\heading}[1]{\section*{\color{accent}#1}}
\newcommand{\slideheading}[1]{\subsection*{\color{accent}#1}}
\newcommand{\cue}[2]{\par\noindent\textbf{#1}\quad #2\par}
\newcommand{\slideguide}[3]{\slideheading{Slide #1 | #2}\cue{Guide}{#3}}
\setlength{\emergencystretch}{1em}
\newcommand{\E}{\mathbb{E}}
\newcommand{\Var}{\operatorname{Var}}
\newcommand{\Cov}{\operatorname{Cov}}
\begin{document}
\begin{center}{\Large\bfseries Joint tests and interpretation}\end{center}
\textbf{Slide route:} Lecture6v2.tex: single coefficients, joint hypotheses, F tests, linear restrictions and confidence sets. Chapter 7.
\heading{Session 8 | A claim about several coefficients}
\textbf{Outcomes:} Distinguish separate from joint tests, calculate a classical F statistic, test a linear combination, and interpret restrictions in context.
\heading{100-minute route}
\begin{tabularx}{\textwidth}{@{}rX@{}}\toprule
0--10 & Partial-coefficient retrieval.\\
10--25 & Individual versus joint hypotheses.\\
25--45 & F-statistic and restricted/unrestricted example.\\
45--55 & Robust versus classical tests and pause.\\
55--75 & Single restrictions on multiple coefficients.\\
75--90 & Confidence sets and model specification.\\
90--100 & Exit ticket and functional-form preview.\\\bottomrule
\end{tabularx}
\heading{Slide-by-slide script}
\slideguide{1}{Title}{Introduce the class as learning how to test a claim about several coefficients at once.}
\slideguide{2}{Roadmap}{Preview individual tests, joint hypotheses, the F statistic, restrictions, and confidence sets.}
\slideguide{3}{Last lesson}{Retrieve the meaning of a partial coefficient and ask why one coefficient test may not answer a policy question.}
\slideguide{4}{Last lesson}{Use the dummy-variable example to remind students that a group comparison may involve more than one coefficient.}
\slideguide{5}{Single coefficients}{Distinguish $H_0:\beta_j=0$ from a joint claim about several $\beta$s.}
\slideguide{6}{Joint hypotheses}{Write $H_0:\beta_2=\beta_3=0$ before calculating. Ask what the restricted model removes.}
\slideguide{7}{Joint hypotheses}{Explain why separate 5 percent tests do not implement one joint 5 percent test.}
\slideguide{8}{Joint hypotheses}{Use the restricted and unrestricted residual sums of squares as the two objects in the F calculation.}
\slideguide{9}{F statistic}{Write $F=((SSR_R-SSR_U)/q)/(SSR_U/(n-k-1))$ and label numerator and denominator.}
\slideguide{10}{F statistic}{With 120 and 100, $q=2$, $n=100$, $k=3$, compute $F=9.6$.}
\slideguide{11}{F distribution}{Explain that the reference distribution depends on the restrictions and denominator degrees of freedom.}
\slideguide{12}{F distribution}{Ask what a large F means: the restrictions worsen fit relative to residual noise, under the maintained assumptions.}
\slideguide{13}{Single coefficients}{Return to a t test and explain that it addresses one restriction.}
\slideguide{14}{Homoskedastic-only F}{State that the residual-sum-of-squares formula is not automatically robust.}
\slideguide{15}{Homoskedastic-only F}{Use the displayed regression output to identify $q$, $k$, and the relevant SSRs.}
\slideguide{16}{Example with class-size variables}{Work the deck's joint test and ask students to state the conclusion without overclaiming each coefficient.}
\slideguide{17}{Example continued}{Compare the joint conclusion with individual p-values. They need not tell the same story.}
\slideguide{18}{Why the F statistic?}{Explain that one statistic aggregates how far the restrictions move the fit.}
\slideguide{19}{Summary}{Have students say when a joint test is more natural than a list of individual tests.}
\slideguide{20}{Single restrictions}{Introduce $H_0:\beta_1=\beta_2$ as one restriction, $\beta_1-\beta_2=0$.}
\slideguide{21}{Single restrictions}{Use $\Var(\hat\beta_1-\hat\beta_2)$ and emphasize the covariance term.}
\slideguide{22}{Rearranging}{Show that rewriting a restriction changes notation, not the hypothesis.}
\slideguide{23}{Direct testing}{Use a software command only after the hand calculation is clear.}
\slideguide{24}{Confidence sets}{Explain that a joint confidence region must account for coefficient covariance.}
\slideguide{25}{Confidence sets}{Contrast a confidence ellipse with the Cartesian product of two individual intervals.}
\slideguide{26}{Confidence sets}{Ask which parameter combinations are ruled out by the region.}
\slideguide{27}{Confidence sets}{Use the interaction preview: a conditional effect is a linear combination that may need its own interval.}
\slideguide{28}{Specification}{Ask which variables belong in a model based on the question and design, not only on significance.}
\slideguide{29}{Specification}{Warn that dropping controls because they are individually insignificant can reintroduce OVB.}
\slideguide{30}{Summary}{Review joint null, restriction count, covariance, and reference distribution.}
\slideguide{31}{Material}{Assign one robust Wald exercise and one equality restriction before the functional-form class.}

\heading{Worked example | Two exclusions}
\cue{Use with slides}{Lecture6v2 slides 7--16: Tests of Joint Hypotheses through the example on $strs$ and $frpmfracs$. Use the SSR calculation when slide 11 introduces the F-statistic, then compare it with the deck's example on slide 16.}
Test $H_0:\beta_2=\beta_3=0$. A restricted regression has residual sum of squares 120; the unrestricted regression has 100. Let $n=100$, unrestricted slope count $k=3$, and restriction count $q=2$.
\[
F=\frac{(SSR_R-SSR_U)/q}{SSR_U/(n-k-1)}
=\frac{(120-100)/2}{100/96}=9.6.
\]
Under homoskedastic normal errors, compare with $F_{2,96}$; its 5\% critical value is about 3.09, so reject.
\cue{What to say}{The joint claim that both coefficients are zero is inconsistent with the data under these assumptions. This does not establish that each individual coefficient is nonzero.}
\cue{Robust caveat}{The residual-sum-of-squares formula is the homoskedastic-only test. Under heteroskedasticity, use the estimated robust covariance matrix in a Wald test; do not relabel this formula robust.}

\newpage
\heading{Worked example | Equality is one restriction}
\cue{Use with slides}{Lecture6v2 slides 22--24: Testing Single Restrictions, Rearranging, and Direct Testing. Use this equality example before the Stata command so students see why the covariance term matters.}
Suppose $\hat\beta_1=.5$, $\hat\beta_2=.2$, their estimated variances are $.04$ and $.01$, and covariance is $.01$. To test $H_0:\beta_1=\beta_2$:
\[
\widehat{\beta_1-\beta_2}=.3,\qquad
SE(\hat\beta_1-\hat\beta_2)=\sqrt{.04+.01-2(.01)}=\sqrt{.03}.
\]
\[
t=.3/\sqrt{.03}\simeq1.732,\quad
95\%\text{ large-sample CI: } .3\pm1.96\sqrt{.03}
=[-.0395,.6395].
\]
Fail to reject equality at 5\%. Each coefficient's separate significance relative to zero cannot answer whether the two coefficients differ.
\cue{Pair prompt}{Why can we not subtract the two standard errors? Because the variance of a difference includes both variances and their covariance.}
\cue{Multiple-testing side example}{If two tests were independent and each used 5\%, the probability of at least one false rejection under the joint null would be $1-.95^2=.0975$. Regression coefficient tests need not be independent, but this shows why separate tests do not implement a joint 5\% test.}
\cue{Software illustration}{With an appropriately fitted Stata model:}
\begin{verbatim}
test x2 x3
test x1 = x2
\end{verbatim}
Use the model's robust covariance estimate when robust inference is required. Confirm variable names before demonstrating.
\cue{Confidence sets}{A joint confidence ellipse incorporates coefficient covariance. It is not obtained by declaring the Cartesian product of individual 95\% intervals a joint 95\% region.}
\cue{Exit answers}{For one restriction, a consistently constructed Wald F equals the corresponding $t^2$. Equality of two coefficients is one restriction. An insignificant individual test does not preclude joint significance.}
\cue{If short / preparation}{Leave repeated output screenshots for reading. Preserve both worked examples. Read Chapter 8 and ask students to sketch a relationship whose slope changes with $X$.}
\textbf{Source:} Lecture6v2.tex; synthetic numerical examples.

\end{document}
